\documentclass[11pt,a4paper]{article}

\usepackage[utf8]{inputenc} % allow utf-8 input
\usepackage[T1]{fontenc}    % use 8-bit T1 fonts
\usepackage{hyperref}       % hyperlinks
\usepackage{url}            % simple URL typesetting
\usepackage{booktabs}       % professional-quality tables
\usepackage{amsfonts}       % blackboard math symbols
\usepackage{nicefrac}       % compact symbols for 1/2, etc.
\usepackage{microtype}      % microtypography
\usepackage{xcolor}         % colors
\usepackage{graphicx}       % graphics
\usepackage[normalem]{ulem} % strikethrough
\title{HTML Text Formatting Examples}

\begin{document}

\maketitle

This is a \textbf{bold (b)} example and right next to it we have a \textbf{strong emphasis (strong)}. Notice that \textbf{strong + bold mixed} looks similar but carries additional semantic meaning.

Here is an \textit{italic (i)} word and an \textit{emphasis (em)} example. Sometimes we combine them like \textit{italic + emphasis together}.

Now let's look at text that appears crossed out: \sout{strikethrough with s} and \sout{deleted with del}. You can also mix them: \sout{double strikethrough (s + del)}.

To highlight insertions or underlines: underlined with u, inserted with ins. A combination could be: underline + insertion together.

Subscript and superscript examples: Water is written as H2O using sub. The mathematical expression x2 + y3 uses sup. They can also be combined: CO2*.

Mixing several: This sentence has \textit{\textbf{strong + emphasis}}, some \textbf{bold + underline}, and a formula like a2 + b3.

\end{document}